\section{Classical comparisons and the remaining lift problem}\label{sec:scope}
\subsection{Fixed degree and positive order}
The direct precedent for fixed degree three with unbounded positive order
is Benvenuti--Farina \cite{BF99}, not merely the familiar stationary
irrational-rotation obstruction. In the normalization reproduced by
Czaja--Jaming--Matolcsi \cite[Example 2, equation (3.1)]{CJM}, their family is
\[
 H_m(z)=\frac1{z-1}-\frac{4(5/2)^{m-2}}{z-2/5}
                       +\frac{3\,5^{m-2}}{z-1/5},\qquad m\ge4.
\]
Its McMillan degree is three. The 1999 lower bound on positive order is
$m$, and the later construction attains $m$ \cite{BF99,CJM}.
Those statements concern one repeated nonnegative matrix and an entire
impulse response. They are not lower bounds against independently
chosen stochastic matrices at every cut of a finite word experiment.

There is an equally direct precedent for the inadequacy of positive
factorization without dynamics. Benvenuti--Farina
\cite[Section 3, Example 4 and Theorem 2]{BF03} give a periodic Hankel
matrix with a nonnegative factorization of inner size three but no
three-dimensional positive realization; their minimality theorem includes
a positive intertwining matrix. Thus neither the need for compatibility
nor its failure in positive factorization is claimed as new here.
For a finite prefix--suffix array, conditioning supplies a positive
factorization at a cut, just as for a finite Hankel truncation. The issue
in this article is one sequence of such factorizations under every
command, allowing unrelated rows and a horizon-dependent machine.
The profile budget proves a necessary inequality even in that enlarged
class. It is not claimed to sharpen the known stationary lower bound
for the family $H_m$.

Accessible spaces, observable quotients and invariant cones are classical
\cite{Heller,BF04,Vidyasagar,MW}. Lemma~\ref{lem:reachable} is the
finite-word version needed to avoid asserting an intertwining identity
on hidden directions that the observations do not determine.
The proof is included to expose that quantifier, not to claim invention
of an accessible-space construction.

\subsection{Static and dynamic extensions}
Yannakakis' factorization theorem relates extension size to nonnegative
rank of a slack matrix \cite{Yannakakis}; see the explicit formulation in
\cite[Theorem 1]{FRT}. The number of faces of an extension supplies the
standard logarithmic lower estimate from projected vertices.
Fiorini--Rothvo\ss--Tiwary \cite[Theorem 2]{FRT}, building on
Ben-Tal--Nemirovski \cite{BTN}, show that a regular $m$-gon has an
$O(\log m)$ extension. Vandaele--Gillis--Glineur \cite{VGG} refine
regular-polygon bounds using slack factorizations.

These theorems minimize a static inequality count. They do not require
the lifted command maps or the observable decoder to extend to stochastic
maps on the whole ambient simplex. Proposition~\ref{prop:slack} shows
that bounded polytope extensions can be normalized without increasing
their facet count; Proposition~\ref{prop:extension-gap} shows why this
normalization alone does not settle dynamical implementation.
The extension test \eqref{eq:dual} is elementary convex separation.
Its role is to make the missing condition exact and checkable for fixed
lift data, not to claim a new general linear-programming theorem.

The substantive quantitative conclusion is therefore divided in two.
For arbitrary hidden realizations, the new potential yields a logarithmic
lower bound, while the resonant construction yields a cube-root upper
bound. For vector-state realizations, the same potential has no
exponential facet-to-vertex conversion, and the order is exactly
$N^{1/3}$. An independent small static lift does not bridge this gap.
Conversely, the map-defect example does not rule out small compatible
lifts that have a different geometry.

\subsection{What is retained, and what is not inferred}
The older volume obstruction, rank-three two-cut example, exact tagged
profile identity, online query constructions and physical saddle results
remain in their original sources and in the supporting archives
\cite{GTF31,GTF32}. In particular the tagged identity still requires
exact final tag recovery and a common suffix interface. Nothing in the
present paper extends it to noisy tags or overlapping supports.
The earlier additive replacement update is a weighted reservoir
construction; its sampling mechanism is credited to that literature
\cite[Appendix A.2]{PBRT}, not used as a new claim here.

The finite proofs above are self-contained apart from the named classical
background identities whose short arguments are also given. The source
record distinguishes the original 1999 reference from the full author
survey and later matching-order proof used to inspect its family. The
original 1999 article and the original Yannakakis article have not been
claimed to receive an exhaustive proof-by-proof priority audit. Their
results are expressly credited; the full accessible comparisons
\cite{BF03,CJM,FRT,VGG} supply the concrete theorem statements used here.
No exhaustive originality clearance is inferred from this bibliography.

The main unresolved issue is the optimal hidden width between logarithmic
and cube-root growth. The vector-state result is not substituted for that
question. The rational-angle bound is explicit but much weaker; its
sharp denominator-controlled rate is likewise not identified.

The autonomous theorem concerns one fixed decoder at all stopping lengths,
not just removal of a clock from the original single-length specification.
All counts continue to price persistent labels rather than read-only
programs, arithmetic, or physical realization of irrational coins.
These are different optimization problems, and the comparisons depend
on preserving their definitions.

The program title identifies this manuscript within the preserved research
series. It does not assert that finite positive-realization arguments
replace independent Fourier, large-deviation, semigroup, or operator-domain
proof obligations elsewhere in that series, nor does it assert a journal
decision. The theorem chain established here is the one displayed in
Theorem~\ref{thm:main} and its supporting lift and profile statements.
